%% ============================================================
%%  Point-by-point response to Reviewer 1 · Round 1
%%  engproc-4523067 · Engineering Proceedings (MDPI) · ICI2ST 2026
%%  Standalone: compile with pdflatex.
%%  AUTOR: revisar los números de sección contra el PDF final antes de enviar.
%% ============================================================
\documentclass[11pt,a4paper]{article}

\usepackage[T1]{fontenc}
\usepackage[utf8]{inputenc}
\usepackage{lmodern}
\usepackage[english]{babel}
\usepackage{amsmath}
\usepackage[left=2.3cm,right=2.3cm,top=2.3cm,bottom=2.3cm]{geometry}
\usepackage{microtype}
\usepackage{enumitem}
\usepackage{xcolor}
\usepackage{mdframed}
\usepackage[hidelinks]{hyperref}

\definecolor{cmtbg}{HTML}{F1F5F9}
\definecolor{cmtrule}{HTML}{94A3B8}
\newenvironment{comment}
  {\begin{mdframed}[backgroundcolor=cmtbg, linecolor=cmtrule, linewidth=0.6pt,
     innertopmargin=5pt, innerbottommargin=5pt, innerleftmargin=7pt,
     innerrightmargin=7pt, skipabove=4pt, skipbelow=6pt]\itshape}
  {\end{mdframed}}
\newcommand{\resp}{\par\noindent\textbf{Response.}\ }
\newcommand{\done}[1]{\par\vspace{2pt}\noindent\textit{Where:} #1\par\vspace{4pt}}
\setlist[itemize]{leftmargin=1.3em, itemsep=2pt, topsep=3pt}
\pagestyle{plain}

\begin{document}

\begin{center}
{\Large\bfseries Response to Reviewer 1}\\[4pt]
{\large \emph{LegalShadow: An Ontology-Guided Digital Shadow Architecture for\\ Semantic Validation of LegalTech Ecosystems}}\\[5pt]
Manuscript ID: engproc-4523067 --- \emph{Engineering Proceedings} (MDPI) --- \today
\end{center}

\vspace{4pt}
\noindent Dear Reviewer,

\vspace{3pt}
\noindent Thank you for a report that was both encouraging and specific. Your structural observations were the most useful part: separating the discussion from the conclusions and giving the policy material a section of its own made the argument easier to follow than it was before, and the abstract is a genuine improvement. We have addressed every point and mark each revised passage in \textcolor{blue}{blue} in the manuscript, with \verb|%>> R1-n| comments in the \LaTeX{} source.

\section*{1. Title}
\begin{comment}
The authors should carefully revise the manuscript title to ensure it accurately reflects the study's main focus and scope.
\end{comment}
\resp Agreed, and this was already acted upon between rounds. The version you reviewed was titled \emph{``\dots Semantic Validation of Judicial Digital Ecosystems''}. The revised title is:
\begin{quote}
\emph{LegalShadow: An Ontology-Guided Digital Shadow Architecture for Semantic Validation of \textbf{LegalTech} Ecosystems}
\end{quote}
The new wording places the paper in the LegalTech research line to which it belongs and matches the governance ontology the instrument validates against, in which the LegalTech domain is one of the ten dimensions. The Academic Editor has reviewed and approved this change. If you consider that the scope should also be signalled in the title, we would be glad to adopt a subtitle such as \emph{``: Evidence from a Judicial Institution''}.

\section*{2. Abstract}
\begin{comment}
The abstract currently does not meet the journal's requirements\dots The authors are required to revise the manuscript to include the study's main objectives, methodology, key findings, and contributions.
\end{comment}
\resp The abstract has been rewritten so that the four elements appear explicitly and in order, and it now sits within the journal's limit (185 words against 206 previously). It states the \emph{objective} as a research question; the \emph{methodology} as the ontology-driven acquisition protocol with SHA-256 provenance, JSON-LD serialisation and the $\langle\Psi,\delta\rangle$ scoring; the \emph{key findings} with all figures preserved (46/100 declarative against $\Psi_{\mathrm{global}}=34.9\%$, 90.4\% for the control, three dimensions at zero); and the \emph{contributions} as the reference architecture, the auditable protocol and the coverage metric.
\done{Abstract, fully rewritten.}

\section*{3. Structure of the Introduction}
\begin{comment}
The introduction section has a fundamental structural issue\dots The authors have used subsections for every paragraph, which is unclear and disrupts the flow.
\end{comment}
\resp You identified something real, although its cause lay slightly elsewhere, and we are grateful for the prompt to look. The Introduction contains no subsections; what produced the effect you describe were two devices. First, the three contributions were set with forced line breaks and bold labels \textbf{(C1) (C2) (C3)}, so they read as three headings. Second, and more importantly, Sections~4 to~6 carried twelve bold run-in headings, one on almost every paragraph, so the paper looked subsectioned when skimmed.

We have converted the contributions into running prose and reduced the run-in headings from twelve to five, keeping only those that mark a genuine change of subject. The Introduction now reads as five continuous paragraphs.
\done{Introduction, final paragraph; Sections~4 and~5 throughout.}

\section*{4. Literature Review}
\begin{comment}
Please replace the `Related Work and Delimitation' section with a `Literature Review' section\dots Also, please carefully add references. I have noted that some statements are missing citations. Please cite legal references and closely related research articles as well.
\end{comment}
\resp The section is now titled \textbf{Literature Review} (Section~2). We retained the delimitation table, since it is what substantiates the novelty claim, but it is now one element of the review rather than its organising principle.

On legal references, we agree this was a real omission in a paper about judicial systems, and we have added a fourth thematic block, \emph{Law, legal informatics and the governance of legal technology}, covering the tradition of formal legal knowledge representation \cite{BenchCapon2012}, the Semantic Web agenda for law \cite{Casanovas2016}, quantitative legal analytics \cite{Katz2017}, and the organisational risks specific to e-justice information systems \cite{Rosa2013}. This block also makes explicit a premise the paper depends on but had not stated: that both traditions presuppose the machine-processable availability that this study measures.
\done{Section~2, new subsection and title.}

\section*{5. Recommended references}
\begin{comment}
I recommend that the authors refer to the following paper\dots \emph{The Risk of Global Environmental Change to Economic Sustainability and Law\dots}; \emph{Integrating Artificial Intelligence into Service Innovation, Business Development, and Legal Compliance: Insights from the Hainan Free Trade Port Era}; \emph{TOW environmental migrants in the international refuge law and human rights\dots}
\end{comment}
\resp We read all three carefully and have incorporated one of them.

\textbf{Included.} Li et al., \emph{Integrating Artificial Intelligence into Service Innovation, Business Development, and Legal Compliance} (\emph{Systems}, 2024), is now cited in Section~2 and again in Section~8. Its treatment of legal risk, data protection and algorithmic accountability as design constraints rather than afterthoughts is directly relevant to the AI-integration dimension of our ontology, which is the dimension with the largest measured gap ($\delta=82.4$), and it supports the argument in Policy Implications that structured disclosure is a low-cost regulatory lever.

\textbf{Respectfully not included.} Cao et al.\ on global environmental change and economic sustainability, and Khaskheli et al.\ on environmental migrants in international refugee law, address subject matter we could not connect to the paper's argument without a citation that would not inform the reader: neither concerns digital twins or shadows, ontology-based conformance, open-data maturity, or the governance of judicial information systems. We followed the journal's own guidance on this point, which asks authors to include recommended references only where they enhance the manuscript. We would of course reconsider if you can indicate the specific passage where you see them contributing --- it is entirely possible that we have missed the connection you had in mind.

We did act on the underlying request, which we took to be that a paper about the legal system should engage legal scholarship: hence the new legal block described in~\S4 above.

\section*{6. Figure 2 and description of tables}
\begin{comment}
Figure 2 is unclear and should be improved for better readability. In addition, the authors should provide a clear description for each table included in the manuscript.
\end{comment}
\resp \textbf{Figure 2} (the JSON-LD fragment) has been re-typeset as a numbered listing with a larger effective size, syntax colouring that distinguishes string values from structural keys, and highlighting of the three fields that carry the argument. Its caption now walks the reader through the object line by line: which lines hold the evidence block, which constitute the anchoring $\Gamma(v)$, and which record the gap.

\textbf{Tables.} Both now have a descriptive paragraph in the running text. For Table~1 we explain what each row and each of the three columns represents and what reading across a row versus down the last column shows. For Table~2 we explain every column, the ordering, the aggregate row, and how the \emph{Root} column alone already separates the three coverage regimes the metric distinguishes.

We also took the opportunity to correct a typographic defect we found in Figure~3(a): the text of two pipeline stages overflowed the boxes containing it. The boxes have been widened and the stage spacing adjusted.
\done{Figure~2 and caption; Section~2 after Table~1; Section~5 before Table~2; Figure~3(a).}

\section*{7. Separating Discussion and Conclusions}
\begin{comment}
The Discussion and Conclusions sections should be written separately. Currently, both are very confusing and difficult to follow.
\end{comment}
\resp Done. \textbf{Section~6 (Discussion)} now retains only interpretation: the contrast between declarative maturity and semantic coverage, the reflexive finding and the measurement bound it implies, external validity with the transfer protocol, and limitations. \textbf{Section~7 (Conclusions)} is new and synthetic, stating three conclusions --- that declarative maturity and machine-actionable governance are distinct quantities, that the deficit is localisable rather than diffuse, and that the instrument is auditable end to end.
\done{Sections~6 and~7.}

\section*{8. New sections after the Conclusion}
\begin{comment}
The authors need to add the following sections after the Conclusion and also revise the Conclusion: add sections Policy Implications, Future Research Directions, and Recommendations.
\end{comment}
\resp All three have been added as Sections~8, 9 and~10, and the Conclusion has been rewritten as described above.

One note on how we did it. The material these sections require already existed in the manuscript, distributed across the old combined section: the prescriptive half of the reflexive finding, the closing list of future work, and the three concrete publication acts. Rather than write new sections alongside it, we \emph{moved} that content into them, so the paper says each thing once. Section~8 (\textbf{Policy Implications}) draws out three implications for how judicial institutions are assessed and how they report. Section~9 (\textbf{Future Research Directions}) sets out four directions, each traced to a specific limitation. Section~10 (\textbf{Recommendations}) gives three publication acts for judicial institutions and one for oversight bodies, with the coverage dimensions each would unlock.

The manuscript grows from 12 to 14 pages as a result. We are happy to compress further if the volume prefers a shorter format.
\done{Sections~8, 9 and~10, all new.}

\vspace{6pt}
\hrulefill
\vspace{4pt}

\noindent No scientific content, numerical result or claim has changed in this revision; every figure in the paper is unchanged in value and remains recomputable from the artefacts deposited at Zenodo (DOI:~10.5281/zenodo.21942315).

\vspace{6pt}
\noindent Thank you again for the care you gave the manuscript.

\vspace{12pt}
\noindent Yours sincerely,\\[14pt]
\textbf{Patricio M. Paccha-Angamarca}, on behalf of all authors\\
\href{mailto:patricio.paccha@epn.edu.ec}{patricio.paccha@epn.edu.ec}

\begin{thebibliography}{9}\footnotesize
\bibitem{BenchCapon2012} Bench-Capon, T.; et al. A history of AI and Law in 50 papers. \emph{Artif. Intell. Law} \textbf{2012}, \emph{20}, 215--319.
\bibitem{Casanovas2016} Casanovas, P.; et al. Semantic Web for the Legal Domain: The next step. \emph{Semantic Web} \textbf{2016}, \emph{7}, 213--227.
\bibitem{Katz2017} Katz, D.M.; et al. A general approach for predicting the behavior of the Supreme Court. \emph{PLoS ONE} \textbf{2017}, \emph{12}, e0174698.
\bibitem{Rosa2013} Rosa, J.; et al. Risk factors in e-justice information systems. \emph{Gov. Inf. Q.} \textbf{2013}, \emph{30}, 241--256.
\end{thebibliography}

\end{document}
